\documentclass[11pt]{article}
\usepackage[margin=1.8cm]{geometry}
\usepackage{graphicx}
\usepackage{booktabs}
\usepackage{array}
\usepackage{longtable}
\usepackage{tabularx}
\usepackage{hyperref}
\hypersetup{hidelinks}
\setlength{\parindent}{0pt}
\setlength{\parskip}{5pt}

\begin{document}

\begin{center}
{\Large\bfseries Supplementary material}\par
\vspace{6pt}
{\large Self-supervised biomedical signal representation learning from intrapartum cardiotocography: stability, ablation, and exploratory neonatal risk enrichment}
\end{center}

\section*{Supplementary Figure S1}
\begin{center}
\includegraphics[width=\textwidth]{figures/Figure_S1_R5_model_selection_ablation.png}
\end{center}
\textbf{AI model selection and sensitivity.} (A) Mean development AUPRC and standard error for four prespecified architectures. The vertical dashed line is the one-standard-error eligibility threshold; p80-d48-l2 mixed was the smallest eligible model. (B) Mask-strategy sensitivity on a fixed development split. (C) Objective ablation including the random-encoder negative control. (D) Fixed-test dimensionality, regularization and missingness sensitivities. Panels B and C were not used to reselect the final architecture; Panel D is exploratory because the fixed test set had prior exposure.

\clearpage
\section*{Supplementary Figure S2}
\begin{center}
\includegraphics[width=\textwidth]{figures/Supplementary_Figure_S2_R5_external_representation.png}
\end{center}
\textbf{External representation transport and domain shift.} (A) Principal-component projection of selected-encoder embeddings from CTU-UHB and CTGDL-FHRMA. (B) Domain-classifier ROC curve (AUROC 0.8811). (C) CTGDL-FHRMA assignments to CTU-fitted SSL phenotype centroids. (D) Distance to the nearest CTU phenotype centroid. This is representation/domain-shift validation, not external neonatal-outcome validation.

\clearpage
\section*{Supplementary Figure S3}
\begin{center}
\includegraphics[width=\textwidth]{figures/Supplementary_Figure_S3_JNU_cross_domain_sensitivity.png}
\end{center}
\textbf{Public JNU-CTG cross-domain pretraining sensitivity.} (A) Five-seed development out-of-fold AUPRC by encoder origin. The dotted line denotes development outcome prevalence. (B) Macro-AUROC across seven CTU expert-morphology probes; classical signal features are a supervised conventional-feature comparator. (C) Paired hierarchical bootstrap differences for JNU pretraining followed by CTU adaptation. JNU transfer improved morphology probes relative to CTU-only and random embeddings, whereas its risk-enrichment interval crossed zero. (D) Exploratory fixed-test AUPRC after the development promotion decision; the dotted line denotes test outcome prevalence. (E) Median AUPRC change for FHR-only, UC-only and 1-Hz inputs relative to full 4-Hz FHR+UC. Directions differed between development and fixed-test analyses. (F) Channel-asymmetric minus mixed JNU masking. The morphology interval was negative and the risk interval crossed zero. Fixed-test panels retain potential selection-induced optimism; no JNU result represents external outcome validation.

\clearpage
\section*{Supplementary Table S1. Fixed cohort split}
\begin{tabular}{lrrr}
\toprule
Split & Records & Events & Prevalence \\
\midrule
Development & 386 & 79 & 0.2047 \\
Fixed test & 166 & 34 & 0.2048 \\
\bottomrule
\end{tabular}

\section*{Supplementary Table S2. Development-only architecture selection}
{\small
\begin{tabularx}{\textwidth}{>{\raggedright\arraybackslash}Xrrrr}
\toprule
Architecture & Mean AUPRC & SE & Parameters & Selected \\
\midrule
p60-d64-l2 mixed cw0.2 & 0.4045 & 0.0826 & 121,532 & No \\
p120-d64-l2 block cw0.2 & 0.3962 & 0.0636 & 126,712 & No \\
p80-d48-l2 mixed cw0.1 & 0.3938 & 0.0897 & 74,768 & Yes \\
p80-d64-l3 channel cw0.2 & 0.3900 & 0.0832 & 172,816 & No \\
\bottomrule
\end{tabularx}
}

The highest mean AUPRC was 0.4045; its one-standard-error threshold was 0.3218. The smallest eligible architecture was selected. Fold-level results and split assignments are provided as CSV files.

\section*{Supplementary Table S3. Development-only objective and mask sensitivities}
\begin{tabular}{lrrr}
\toprule
Objective & AUPRC & AUROC & Brier \\
\midrule
Random encoder & 0.3333 & 0.6883 & 0.2385 \\
Reconstruction only & 0.3174 & 0.6656 & 0.2408 \\
NT-Xent only & 0.2991 & 0.6695 & 0.2433 \\
Reconstruction + NT-Xent & 0.3106 & 0.6760 & 0.2448 \\
Combined, no jitter & 0.3100 & 0.6786 & 0.2453 \\
\bottomrule
\end{tabular}

\vspace{8pt}
\begin{tabular}{lrrr}
\toprule
Mask strategy & AUPRC & AUROC & Brier \\
\midrule
Random & 0.3461 & 0.7104 & 0.2300 \\
Block & 0.3312 & 0.7039 & 0.2372 \\
Channel & 0.3296 & 0.6818 & 0.2352 \\
Mixed & 0.3106 & 0.6760 & 0.2448 \\
\bottomrule
\end{tabular}

All results used the same fixed 289/97 development split. They are component sensitivities, not test-set selection results.

\section*{Supplementary Table S4. Five-seed stability and paired AUPRC differences}
\begin{tabular}{rrrrrr}
\toprule
Seed & AUPRC & AUROC & Brier & $\Delta$AUPRC & 95\% CI \\
\midrule
20260521 & 0.3929 & 0.5996 & 0.2636 & 0.0164 & $-0.0596$, 0.0763 \\
20260522 & 0.3935 & 0.5867 & 0.2628 & 0.0170 & $-0.0623$, 0.0864 \\
20260523 & 0.4042 & 0.6058 & 0.2628 & 0.0277 & $-0.0597$, 0.1013 \\
20260524 & 0.4021 & 0.5934 & 0.2623 & 0.0257 & $-0.0551$, 0.0886 \\
20260525 & 0.3916 & 0.5954 & 0.2643 & 0.0151 & $-0.0665$, 0.0880 \\
\bottomrule
\end{tabular}

The signal baseline was fixed across seeds (AUPRC 0.3765, AUROC 0.6428, Brier 0.2482). Differences are reduced SSL minus signal baseline using 2,000 paired bootstrap resamples.

\section*{Supplementary Table S5. Dimensionality, regularization and missingness sensitivities}
\begin{tabular}{lrrrr}
\toprule
Model & Predictors & EPV & AUPRC & AUROC \\
\midrule
Signal 10 & 10 & 7.90 & 0.3765 & 0.6428 \\
Signal 10 + SSL PCA5, L2 & 15 & 5.27 & 0.3929 & 0.5996 \\
Signal 10 + SSL PCA10, L2 & 20 & 3.95 & 0.4229 & 0.6085 \\
Signal 10 + SSL PCA20, L2 & 30 & 2.63 & 0.4210 & 0.6179 \\
Signal 10 + SSL PCA5, elastic net & 15 & 5.27 & 0.3918 & 0.6005 \\
Signal 12 + SSL PCA5, L2 & 17 & 4.65 & 0.3967 & 0.6217 \\
\bottomrule
\end{tabular}

EPV, events per predictor, uses 79 development events. PCA10/20 and missingness results were not used to replace the prespecified PCA5 primary model.

\section*{Supplementary Table S6. Phenotype incremental sensitivity}
\begin{tabular}{llrrr}
\toprule
$k$ & Model & AUPRC & AUROC & Predictors \\
\midrule
3 & Reduced SSL & 0.3929 & 0.5996 & 15 \\
3 & + signal phenotype & 0.3796 & 0.5762 & 17 \\
3 & + SSL phenotype & 0.3804 & 0.6005 & 17 \\
4 & Reduced SSL & 0.3929 & 0.5996 & 15 \\
4 & + signal phenotype & 0.3978 & 0.6047 & 18 \\
4 & + SSL phenotype & 0.3511 & 0.5858 & 18 \\
4 & + both phenotypes & 0.3457 & 0.5831 & 21 \\
5 & Reduced SSL & 0.3929 & 0.5996 & 15 \\
5 & + signal phenotype & 0.3944 & 0.5960 & 19 \\
5 & + SSL phenotype & 0.3736 & 0.5882 & 19 \\
\bottomrule
\end{tabular}

Complete $k=3$, 4 and 5 rows, including both-phenotype models, are supplied as CSV. Phenotypes are secondary interpretation variables.

\section*{Supplementary Table S7. Endpoint sensitivity}
\begin{tabular}{llrrrr}
\toprule
Outcome & Model & Dev events & Test events & AUPRC & AUROC \\
\midrule
Composite & Signal & 79 & 34 & 0.3765 & 0.6428 \\
Composite & Reduced SSL & 79 & 34 & 0.3929 & 0.5996 \\
pH $<$7.15 & Signal & 72 & 33 & 0.4161 & 0.6437 \\
pH $<$7.15 & Reduced SSL & 72 & 33 & 0.4044 & 0.6138 \\
Apgar5 $<$7 & Signal & 15 & 4 & 0.0593 & 0.4599 \\
Apgar5 $<$7 & Reduced SSL & 15 & 4 & 0.0591 & 0.6173 \\
pH $<$7.10 & Signal & 37 & 19 & 0.4507 & 0.7479 \\
pH $<$7.10 & Reduced SSL & 37 & 19 & 0.4588 & 0.7429 \\
pH $<$7.05 & Signal & 25 & 15 & 0.5721 & 0.8283 \\
pH $<$7.05 & Reduced SSL & 25 & 15 & 0.5382 & 0.7991 \\
\bottomrule
\end{tabular}

The Apgar-only analysis has fewer than 10 test events and is reported only as an unstable exploratory estimate.

\section*{Supplementary Table S8. Calibration metrics}
\begin{tabular}{llrrrr}
\toprule
Model & Calibration & Brier & ECE & Intercept & Slope \\
\midrule
Signal & Raw & 0.2482 & 0.2874 & $-1.4267$ & 0.6518 \\
Signal & Platt & 0.1543 & 0.0386 & $-0.2404$ & 0.9006 \\
Signal & Isotonic & 0.1581 & 0.0588 & $-0.5425$ & 0.6506 \\
Reduced SSL & Raw & 0.2636 & 0.2747 & $-1.3920$ & 0.3860 \\
Reduced SSL & Platt & 0.1579 & 0.0915 & $-0.5638$ & 0.6440 \\
Reduced SSL & Isotonic & 0.1694 & 0.0878 & $-0.8558$ & 0.4171 \\
\bottomrule
\end{tabular}

\section*{Supplementary Table S9. Historical exploratory result}
\begin{tabularx}{\textwidth}{p{4.2cm}rrX}
\toprule
Historical model & AUPRC & AUROC & Status \\
\midrule
Signal + 64 SSL dimensions + signal phenotype & 0.4016 & 0.6168 & Exploratory only. The fixed test set contributed to configuration ranking, so this estimate is not independent validation and is excluded from revised model selection and primary reporting. \\
\bottomrule
\end{tabularx}

\section*{Supplementary Table S10. Selected encoder and computing environment}
\begin{tabular}{ll@{\qquad}ll}
\toprule
Parameter & Value & Parameter & Value \\
\midrule
Patch length & 80 & Model dimension & 48 \\
Layers & 2 & Attention heads & 4 \\
Mask strategy & Mixed & Mask fraction & 0.25 \\
Reconstruction weight & 1.0 & Contrastive weight & 0.1 \\
Temperature & 0.2 & Jitter SD & 0.02 \\
Epochs & 4 & Batch size & 64 \\
Learning rate & 0.001 & Weight decay & 0.0001 \\
Parameters & 74,768 & Training windows & 5,301 \\
Framework & PyTorch 2.12.0+cu130 & Device & RTX 5080 Laptop GPU \\
\bottomrule
\end{tabular}

\section*{Supplementary Table S11. External representation analysis}
\begin{tabular}{lr}
\toprule
Metric & Value \\
\midrule
CTU-UHB records & 552 \\
CTGDL-FHRMA records & 135 \\
Embedding dimension & 48 \\
Domain-classifier AUROC & 0.8811 \\
External outcome validation & No \\
Representation/domain-shift analysis & Yes \\
\bottomrule
\end{tabular}

\section*{Supplementary Tables S12--S13. Reproducibility audit}

Supplementary Table S12 provides SHA256 hashes for locked inputs. Supplementary Table S13 provides the 23-item output cross-audit, including fold-membership checks, recomputation of AUPRC, AUROC and Brier scores from record-level probability tables, seed completeness, endpoint event counts and figure-file checks. Both are supplied as machine-readable CSV files. The audit status was PASS.

\section*{Supplementary Table S14. Encoder-origin risk stability}
{\small
\begin{tabularx}{\textwidth}{>{\raggedright\arraybackslash}Xrrrr}
\toprule
Encoder origin & Dev AUPRC & Dev AUROC & Fixed-test AUPRC & Fixed-test AUROC \\
\midrule
Random initialization & 0.3611 & 0.6584 & 0.3840 & 0.5978 \\
CTU-development-only SSL & 0.3797 & 0.6691 & 0.4017 & 0.5945 \\
JNU-only SSL & 0.3658 & 0.6702 & 0.4012 & 0.6092 \\
JNU pretraining $\rightarrow$ CTU adaptation & 0.3682 & 0.6720 & 0.3420 & 0.5951 \\
JNU channel-asymmetric sensitivity & 0.3802 & 0.6475 & -- & -- \\
\bottomrule
\end{tabularx}
}

Values are medians across seeds 20260521--20260525. Development results are three-fold out-of-fold estimates. Fixed-test results are exploratory because the test set had historical exposure. The channel-asymmetric route was development-only.

\section*{Supplementary Table S15. Frozen-encoder morphology probes}
\begin{tabular}{lrr}
\toprule
Representation & Macro-AUPRC & Macro-AUROC \\
\midrule
Classical signal features & 0.5666 & 0.7357 \\
Random encoder & 0.4543 & 0.6211 \\
CTU-development-only SSL & 0.4789 & 0.6392 \\
JNU-only SSL & 0.5155 & 0.6738 \\
JNU pretraining $\rightarrow$ CTU adaptation & 0.5076 & 0.6802 \\
JNU channel-asymmetric sensitivity & 0.4460 & 0.6005 \\
\bottomrule
\end{tabular}

Values are medians across five seeds and seven expert-event tasks. Probes were fitted and evaluated only within record-grouped CTU development folds.

\section*{Supplementary Table S16. JNU paired hierarchical bootstrap comparisons}
{\small
\begin{tabularx}{\textwidth}{>{\raggedright\arraybackslash}p{2.2cm}>{\raggedright\arraybackslash}p{1.9cm}>{\raggedright\arraybackslash}Xr>{\raggedright\arraybackslash}p{2.4cm}r}
\toprule
Data role & Metric & Contrast & Difference & 95\% CI & Positive seeds \\
\midrule
Development & Risk AUPRC & JNU$\rightarrow$CTU vs CTU-only & 0.0030 & $-0.0304$, 0.0316 & 3/5 \\
Development & Morphology AUROC & JNU$\rightarrow$CTU vs CTU-only & 0.0401 & 0.0155, 0.0647 & 5/5 \\
Development & Morphology AUROC & JNU$\rightarrow$CTU vs random & 0.0611 & 0.0344, 0.0884 & 5/5 \\
Development & Morphology AUROC & JNU$\rightarrow$CTU vs signal & $-0.0555$ & $-0.0924$, $-0.0200$ & 0/5 \\
Development & Risk AUPRC & Asymmetric vs mixed JNU & 0.0098 & $-0.0290$, 0.0444 & 4/5 \\
Development & Morphology AUROC & Asymmetric vs mixed JNU & $-0.0787$ & $-0.1006$, $-0.0552$ & 0/5 \\
Fixed test & Risk AUPRC & JNU$\rightarrow$CTU vs CTU-only & $-0.0403$ & $-0.0766$, 0.0016 & 0/5 \\
Fixed test & Risk AUPRC & JNU-only vs CTU-only & 0.0005 & $-0.0204$, 0.0240 & 2/5 \\
\bottomrule
\end{tabularx}
}

Each interval used 2,000 paired hierarchical bootstrap resamples. Fixed-test rows are exploratory and were not used for promotion.

\section*{Supplementary Table S17. Channel and sampling-rate sensitivity}
\begin{tabular}{lllrr}
\toprule
Data role & Encoder & Input & Median AUPRC & $\Delta$ vs full \\
\midrule
Development & CTU-only & FHR-only & 0.3645 & $-0.0153$ \\
Development & CTU-only & UC-only & 0.3600 & $-0.0197$ \\
Development & CTU-only & 1 Hz & 0.3676 & $-0.0122$ \\
Development & JNU$\rightarrow$CTU & FHR-only & 0.3405 & $-0.0276$ \\
Development & JNU$\rightarrow$CTU & UC-only & 0.3789 & 0.0107 \\
Development & JNU$\rightarrow$CTU & 1 Hz & 0.3660 & $-0.0022$ \\
Fixed test & CTU-only & FHR-only & 0.4113 & 0.0096 \\
Fixed test & CTU-only & UC-only & 0.3601 & $-0.0417$ \\
Fixed test & CTU-only & 1 Hz & 0.3910 & $-0.0107$ \\
Fixed test & JNU$\rightarrow$CTU & FHR-only & 0.3961 & 0.0541 \\
Fixed test & JNU$\rightarrow$CTU & UC-only & 0.3521 & 0.0101 \\
Fixed test & JNU$\rightarrow$CTU & 1 Hz & 0.3323 & $-0.0097$ \\
\bottomrule
\end{tabular}

Values are five-seed medians. The direction reversal across data roles precludes selecting an input variant from these sensitivities.

\section*{Supplementary Table S18. JNU-CTG data and leakage audit}
\begin{tabularx}{\textwidth}{p{5.1cm}X}
\toprule
Audit item & Result \\
\midrule
Archive identifier & Zenodo 21800730; CC BY 4.0 \\
Archive MD5 & \texttt{ac1cfcba2f1b3336596544211d776771}; matched \\
Waveform readability & 20,769/20,769 records readable (100\%) \\
Patient grouping & 12,606 unique patient groups retained \\
SSL windows & 62,307 non-overlapping 10-minute windows; FHR+UC; 4 Hz \\
JNU labels used in SSL & None; Apgar, asphyxia, FIGO, demographics and other clinical labels withheld \\
CTU fold isolation & Fold-specific adaptation used development-training records only \\
Fixed-test access & Opened only after the development promotion decision; historical exposure disclosed \\
Promotion decision & Supplement-only because the risk-enrichment gate failed \\
External outcome validation & Not performed \\
\bottomrule
\end{tabularx}

\end{document}
